\documentclass[12pt,a4paper]{article}

\usepackage[utf8]{inputenc}
\usepackage[T1]{fontenc}
\usepackage[slovak]{babel}
\usepackage{geometry}
\usepackage{longtable}
\usepackage{hyperref}
\usepackage{enumitem}
\usepackage{listings}

\geometry{margin=2.5cm}
\hypersetup{
  colorlinks=true,
  linkcolor=black,
  urlcolor=blue
}
\setlist[itemize]{noitemsep, topsep=0.3em}
\renewcommand{\arraystretch}{1.2}

\lstdefinestyle{commandstyle}{
  basicstyle=\ttfamily\small,
  breaklines=true,
  frame=single,
  columns=fullflexible,
  keepspaces=true
}
\lstset{style=commandstyle}

\title{Technická dokumentácia projektu}
\author{Daniela Paluchová}
\date{\today}

\begin{document}

\maketitle
\tableofcontents
\newpage

\section{Úvod}

Projekt je webová aplikácia pre projektové riadenie, plánovanie sprintov a analytickú podporu tímovej práce. Podporuje správu projektov, tímov, úloh, epikov, PERT analýzu, RACI maticu, optimalizáciu plánovania a evidenciu používateľskej aktivity.

Riešenie je postavené ako klient-server systém. Frontend tvorí jednostránková aplikácia vo frameworku Quasar/Vue, backend je REST API vo Flasku a dáta sú uložené v PostgreSQL databáze.

\section{Architektúra systému}

Aplikácia je rozdelená na tri vrstvy:

\begin{itemize}
  \item \textbf{Frontend} -- používateľské rozhranie, routing, lokálny stav a volania API.
  \item \textbf{Backend} -- autentifikácia, autorizácia, doménová logika, plánovanie a REST endpointy.
  \item \textbf{Databáza} -- perzistencia používateľov, projektov, úloh, sprintov, epikov a logov.
\end{itemize}

Frontend komunikuje s backendom cez JSON REST API. Chránené požiadavky používajú JWT token v hlavičke \texttt{Authorization}. API má predvolenú adresu \texttt{http://localhost:5000/api}, ktorú možno zmeniť cez premennú \texttt{VITE\_API\_BASE\_URL}.

\section{Technologický stack}

\begin{longtable}{|p{4cm}|p{9.5cm}|}
\hline
\textbf{Oblasť} & \textbf{Použité technológie} \\
\hline
Frontend & Quasar 2, Vue 3, TypeScript, Vue Router, Pinia, Axios, Chart.js, date-fns. \\
\hline
Backend & Flask 3, Flask-SQLAlchemy, SQLAlchemy 2, Flask-JWT-Extended, Flask-CORS, Flask-Migrate. \\
\hline
Databáza & PostgreSQL cez ovládač \texttt{psycopg2-binary}. \\
\hline
Bezpečnosť & JWT tokeny a bcrypt hashovanie hesiel. \\
\hline
Kvalita kódu & ESLint, Prettier, TypeScript strict režim, vite-plugin-checker. \\
\hline
\end{longtable}

\section{Frontend}

Frontendová časť sa nachádza v adresári \texttt{src/}. Konfiguráciu zabezpečuje \texttt{quasar.config.ts}, kde sú nastavené boot súbory \texttt{auth} a \texttt{axios}, globálne štýly, Quasar pluginy \texttt{Notify} a \texttt{Dialog}, router mód \texttt{hash} a TypeScript strict režim.

\subsection{Štruktúra}

\begin{longtable}{|p{4.5cm}|p{9cm}|}
\hline
\textbf{Cesta} & \textbf{Účel} \\
\hline
\texttt{src/pages/} & Stránky aplikácie: projekty, epiky, tím, workload, PERT/RACI, smart sprint planning, profil a admin logy. \\
\hline
\texttt{src/components/} & Znovupoužiteľné komponenty, napríklad epic board, dialógy a optimalizačné návrhy. \\
\hline
\texttt{src/stores/} & Pinia stores pre autentifikáciu, projekty, epiky, tím a plánovanie. \\
\hline
\texttt{src/router/} & Definícia rout a navigačných guardov. \\
\hline
\texttt{src/services/} & API klienti pre komunikáciu s backendom. \\
\hline
\end{longtable}

\subsection{Routing a stav}

Routy sú definované v \texttt{src/router/routes.ts}. Verejné routy sú \texttt{/login} a \texttt{/register}; hlavná časť aplikácie používa layout \texttt{MainLayout.vue}. Vybrané routy sú chránené meta atribútmi \texttt{requiresAuth}, \texttt{requiresManager} a \texttt{requiresAdmin}.

Autentifikačný store \texttt{src/stores/auth-store.ts} drží aktuálneho používateľa, token, stav načítania a chyby. Token a používateľ sú uložené v \texttt{localStorage}. Store poskytuje akcie pre prihlásenie, registráciu, odhlásenie, úpravu profilu a zmenu hesla.

\subsection{API klient}

Hlavný klient je v \texttt{src/services/api.ts}. Je založený na Axios inštancii s timeoutom 10 sekúnd. Request interceptor pridáva Bearer token mimo verejných endpointov \texttt{/auth/login} a \texttt{/auth/register}. Response interceptor spracúva najmä chyby 401, 403 a 500; pri neautorizovanom prístupe vyčistí lokálnu autentifikáciu.

\section{Backend}

Backendová časť je v adresári \texttt{backend/}. Vstupným súborom je \texttt{backend/run.py}, ktorý vytvára aplikáciu cez \texttt{create\_app()} zo súboru \texttt{backend/app/\_\_init\_\_.py}. Predvolené spustenie používa host \texttt{0.0.0.0}, port \texttt{5000} a debug režim podľa \texttt{FLASK\_ENV}.

Aplikačná továreň inicializuje SQLAlchemy, migrácie, JWT, CORS pravidlá, REST blueprinty a endpoint \texttt{/api/health}. Pri štarte sa volá \texttt{db.create\_all()} a doplnkové automatické migrácie pre vybrané stĺpce epikov.

\subsection{Konfigurácia}

Backend používa environment premenné načítané cez \texttt{python-dotenv}. Dôležité premenné sú \texttt{DATABASE\_URL}, \texttt{SECRET\_KEY}, \texttt{JWT\_SECRET\_KEY}, \texttt{JWT\_ACCESS\_TOKEN\_EXPIRES}, \texttt{CORS\_ORIGINS}, \texttt{HOST}, \texttt{PORT} a \texttt{FLASK\_ENV}. Dokumentácia neuvádza lokálne tajné hodnoty ani obsah súboru \texttt{.env}.

\subsection{API prehľad}

\begin{longtable}{|p{4.5cm}|p{9cm}|}
\hline
\textbf{Prefix} & \textbf{Účel} \\
\hline
\texttt{/api/auth} & Registrácia, prihlásenie, profil, zmena hesla a odhlásenie. \\
\hline
\texttt{/api/projects} & Projekty, sprinty, členovia projektu, kritická cesta, konfigurácia a optimalizačné logy. \\
\hline
\texttt{/api/tasks} & CRUD úloh a hromadné aktualizácie. \\
\hline
\texttt{/api/teams} & Správa členov tímu. \\
\hline
\texttt{/api/projects/\{id\}/epics} & Správa epikov projektu a kritická cesta epikov. \\
\hline
\texttt{/api/projects/\{id\}/smart-sprint-planning} & Návrh inteligentného sprint plánu. \\
\hline
\texttt{/api/projects/\{id\}/analyze-pert-raci} & Kombinovaná PERT/RACI analýza. \\
\hline
\texttt{/api/raci-weights} & Konfigurácia váh RACI výpočtov. \\
\hline
\texttt{/api/activity-logs} & Zápis aktivity používateľov a admin prehľady. \\
\hline
\end{longtable}

\subsection{Autentifikácia}

Backend používa JWT access tokeny. Heslá sú ukladané ako bcrypt hash. Autorizácia pracuje s rolami \texttt{admin}, \texttt{manager}, \texttt{developer} a \texttt{viewer}; pomocné autorizačné funkcie sú v \texttt{backend/app/utils/auth.py}.

\section{Dátový model}

Modely sú definované v \texttt{backend/app/models/}. Najdôležitejšie entity:

\begin{itemize}
  \item \texttt{User} -- používateľský účet s emailom, heslom a rolou.
  \item \texttt{Project} -- projekt s nastaveniami, stavom, termínmi, PERT/RACI konfiguráciou a sprint parametrami.
  \item \texttt{Task} -- úloha s prioritou, stavom, story pointmi, PERT odhadmi, RACI rolami, závislosťami a rizikom.
  \item \texttt{Epic} -- väčší projektový celok s vlastníkom, prioritou, termínmi, PERT odhadmi a závislosťami.
  \item \texttt{Sprint} -- časovo ohraničený plán s kapacitou, story pointmi a velocity.
  \item \texttt{TeamMember} -- člen tímu so zručnosťami, kapacitou, vyťaženosťou a dostupnosťou.
  \item \texttt{ActivityLog}, \texttt{OptimizationLog}, \texttt{TaskHistory} -- audit, optimalizácie a história práce.
\end{itemize}

Projekt má väzby na sprinty, úlohy, epiky a projektové roly. Úloha môže patriť projektu, sprintu aj epiku. Epic môže obsahovať viac úloh a môže mať vlastníka z členov tímu.

\section{Doménové služby}

Služby v \texttt{backend/app/services/} oddeľujú plánovaciu a analytickú logiku od rout. Patria sem najmä:

\begin{itemize}
  \item \texttt{smart\_sprint\_planner.py} -- návrh sprint plánu,
  \item \texttt{pert\_raci\_analyzer.py} -- PERT/RACI analýza,
  \item \texttt{risk\_analyzer.py} -- vyhodnotenie rizík,
  \item \texttt{dependency\_optimizer.py} -- optimalizácia závislostí,
  \item \texttt{bottleneck\_analyzer.py} -- hľadanie úzkych miest,
  \item \texttt{task\_splitter.py} a \texttt{task\_merger.py} -- rozdelenie a spájanie úloh.
\end{itemize}

\section{Lokálne spustenie}

Lokálne sa aplikácia spúšťa ako dve samostatné časti: Quasar frontend a Flask backend. Frontendová časť je Quasar aplikácia, preto sa vo vývoji spúšťa cez Quasar CLI. V súbore \texttt{package.json} je príkaz \texttt{npm run dev} namapovaný na \texttt{quasar dev}.

\begin{lstlisting}
npm install
quasar dev
\end{lstlisting}

Alternatívne je možné použiť npm skript:

\begin{lstlisting}
npm run dev
\end{lstlisting}

Príkaz \texttt{quasar dev} spustí vývojový server Quasar/Vite a otvorí aplikáciu v prehliadači podľa nastavenia v \texttt{quasar.config.ts}. Produkčný build frontendu sa vytvorí príkazom:

\begin{lstlisting}
npm run build
npm run lint
\end{lstlisting}

Backend sa spúšťa zo samostatného adresára \texttt{backend/}. Pred spustením je potrebné mať nastavené environment premenné, najmä \texttt{DATABASE\_URL}, \texttt{SECRET\_KEY} a \texttt{JWT\_SECRET\_KEY}.

\begin{lstlisting}
cd backend
python -m venv .venv
.venv\Scripts\activate
pip install -r requirements.txt
python run.py
\end{lstlisting}

Lokálna databáza je PostgreSQL databáza dostupná cez premennú \texttt{DATABASE\_URL}. Po štarte backend vytvorí tabuľky podľa SQLAlchemy modelov.

\section{Nasadenie}

Frontend je nasadený na platforme Vercel. Konfigurácia je v súbore \texttt{vercel.json}: build príkaz je \texttt{npm run build}, výstupný adresár je \texttt{dist/spa} a všetky SPA routy sú prepísané na \texttt{/index.html}. Verejná adresa aplikácie je:

\begin{lstlisting}
https://diplomovka-implementacia13.vercel.app/
\end{lstlisting}

Backendová časť je nasadená na platforme Render ako Flask REST API. Frontend sa na backend pripája cez premennú \texttt{VITE\_API\_BASE\_URL}, ktorá v produkcii obsahuje URL Render služby. Databázová vrstva používa PostgreSQL databázu hostovanú na platforme Neon a backend sa na ňu pripája cez premennú \texttt{DATABASE\_URL}. Citlivé hodnoty, ako databázové heslá a JWT tajné kľúče, sa ukladajú iba v environment premenných nasadzovacieho prostredia.

\section{Kvalita, testovanie a bezpečnosť}

Frontend používa ESLint, Prettier a TypeScript strict režim. Backend deklaruje \texttt{pytest} a \texttt{pytest-flask}, no projekt zatiaľ nemá plnohodnotnú testovaciu konfiguráciu; \texttt{npm test} je iba placeholder.

Bezpečnosť je založená na JWT a bcrypt. Token sa na klientovi ukladá v \texttt{localStorage}, čo je jednoduché riešenie vhodné pre prototyp, ale pri produkcii treba zohľadniť riziká XSS. CORS je vo vývojovej konfigurácii voľnejší, preto je vhodné v produkcii obmedziť povolené originy.

\section{Zhrnutie}

Projekt kombinuje Quasar/Vue frontend, Flask REST API a PostgreSQL databázu. Hlavnou funkcionalitou je projektové plánovanie podporené PERT odhadmi, RACI zodpovednosťami, kapacitami členov tímu a optimalizáciou sprintov.

\end{document}
